% Adopted Lawvere continuation; working spec fragment preserved separately.
\subsubsection{The abstract-set base and the Boolean-valued correspondence}

Lawvere's Nice 1970 article, published in 1971, announces on p.~333
that its Boolean-valued models satisfy choice and correspond to complete
Boolean algebras in the abstract-set base \cite[p.~333]{Lawvere1970}.
Here ``Boolean'' retains the source's conjunction of classical logic and
support splitting, not just the modern classical-logic condition.
The following constructions are independent editorial proofs of that
announcement at an explicitly specified morphism scope. They do not turn
the six-page source into a source-native detailed proof.

\begin{hblemma}{thm:p10-lawvere-bvm-base}{proved=annex:p10.lawvere-bvm.base}{Bounded well-ordering in the source's base}
Let \(\mathcal S\) have the abstract-set package of
Definition~\ref{def:p10-lawvere-sheaves-bvm}. Then \(\mathcal S\)
is nondegenerate, \(1\) is a generator, every epimorphism splits,
and every object admits an internal well-order. The latter construction
uses power objects and a choice function on nonempty subobjects of that
object; it does not require a cumulative universe, Replacement, an
external ordinal, or a comparison with all models of ZFC.
\end{hblemma}

\begin{hbdefinition}{def:p10-lawvere-bvm-matrices}{Complete Boolean algebras and their matrices}
An \(\mathcal S\)-complete Boolean algebra \(B\) is a Boolean
algebra object with \(\bigvee:\mathcal P_{\mathcal S}B\to B\)
satisfying the least-upper-bound property internally. Meets are obtained
by complementation. A morphism preserves Boolean operations and all
\(\mathcal S\)-indexed joins.

The category \(\operatorname{Mat}(B)\) has objects \((I,b,e)\),
where \(I\in\mathcal S\), \(b:I\to B\), and
\(e:I\times I\to B\) satisfy
\[
 e_{ii}=b_i,\quad e_{ij}\le b_i\wedge b_j,\quad
 e_{ij}=e_{ji},\quad e_{ij}\wedge e_{jk}\le e_{ik}.
\]
A map \((I,b,e)\to(J,c,f)\) is a matrix \(r:I\times J\to B\)
with
\[
\begin{gathered}
 r_{ij}\le b_i\wedge c_j,\qquad \bigvee_j r_{ij}=b_i,\qquad
 r_{ij}\wedge r_{ik}\le f_{jk},\\
 e_{ii'}\wedge r_{ij}\le r_{i'j},\qquad
 r_{ij}\wedge f_{jj'}\le r_{ij'}.
\end{gathered}
\]
These are saturated, total, single-valued relations, not functions between
chosen representatives. Identity is \(e\); composition is
\((sr)_{ik}=\bigvee_j r_{ij}\wedge s_{jk}\).
All displayed indexing and inequalities are internal to \(\mathcal S\).
\end{hbdefinition}

\begin{hbtheorem}{thm:p10-lawvere-bvm-presentation}{proved=annex:p10.lawvere-bvm.presentation}{The actual Boolean-algebra presentation of a model}
For a source Boolean-valued model \(p:\mathcal E\to\mathcal S\),
put \(B_{\mathcal E}=p_*T_{\mathcal E}\). This is an
\(\mathcal S\)-complete Boolean algebra, and there is an equivalence
over \(\mathcal S\)
\[
             \mathcal E\simeq\operatorname{Mat}(B_{\mathcal E}).
\]
Each object is represented by an actual effective quotient of a subobject
of \(p^*I\), for an object \(I\) of the base. The functors on maps,
inverse comparison maps, and their triangle identities are supplied in
the annex. No global points of \(\mathcal E\) are assumed sufficient.
\end{hbtheorem}

\begin{hbtheorem}{thm:p10-lawvere-bvm-realization}{proved=annex:p10.lawvere-bvm.realization}{Realization and choice}
For every \(\mathcal S\)-complete Boolean algebra \(B\), let
\(\operatorname{Sh}_{\mathcal S}(B)\) be internal sheaves on the
poset \(B\), where a sieve on \(b\) covers precisely when its join
is \(b\). Its geometric base map \(q\) has
\(q^*=a\Delta\), and
\[
 q_*T\cong B,\qquad
 \operatorname{Sh}_{\mathcal S}(B)\simeq\operatorname{Mat}(B).
\]
It is classical and every epimorphism splits. Therefore it is a Boolean
topos in the source's stronger sense and satisfies \((\Pi)\) and choice.
The same splitting construction applies to every \(\mathcal E\) in
the preceding theorem. It uses a well-order of a base indexing object
and Boolean disjointification, not a choice of a sheaf-global witness.
The one-element Boolean algebra is allowed and realizes the degenerate
topos; excluding it on one side requires excluding it on the other.
\end{hbtheorem}

\begin{hbdefinition}{def:p10-lawvere-bvm-cell-regime}{The specified morphism and 2-cell regimes}
Write \(\mathsf{BVM}^{\leftarrow}(\mathcal S)\) for the
inverse-image-directed bicategory: a 1-cell from \(\mathcal E\) to
\(\mathcal F\) is a finite-limit-preserving left adjoint
\(L:\mathcal E\to\mathcal F\), with its right adjoint and a specified
base isomorphism \(Lp^*\cong q^*\). A 2-cell is a natural
transformation of these left adjoints compatible with that isomorphism.
Composition has the usual composite base isomorphism.
In the conventional geometric direction a 1-cell
\(\mathcal E\to\mathcal F\) instead has inverse image
\(\mathcal F\to\mathcal E\). We keep this opposite direction
separate. The article names a bicategory on p.~333 without printing its
1-cell or 2-cell definitions; the notation here specifies a reconstruction,
not a quotation of an unprinted convention.
\end{hbdefinition}

\begin{hbtheorem}{thm:p10-lawvere-bvm-correspondence}{proved=annex:p10.lawvere-bvm.correspondence}{Maps, adjunctions, and coherence of the correspondence}
There are biequivalences, with the algebra categories regarded as locally
discrete bicategories,
\[
 \mathsf{BVM}^{\leftarrow}(\mathcal S)
       \simeq\mathsf{CBA}(\mathcal S),\qquad
 \mathsf{BVM}_{\mathrm{geom}}(\mathcal S)
       \simeq\mathsf{CBA}(\mathcal S)^{\mathrm{op}}.
\]
A complete homomorphism \(\alpha:B\to C\) acts on Boolean matrices
entry by entry. Its right adjoint on sheaves is the actual functor
\(R_\alpha(Y)(b)=Y(\alpha b)\). Conversely a 1-cell induces its
homomorphism by transport of subobjects of base-indexed objects.
The adjunction bijections, units and counits, reconstruction isomorphisms,
uniqueness of compatible 2-cells, and composition coherence are established
in the annex. Such inverse images are geometric; they need not preserve
power objects or be logical morphisms.
\end{hbtheorem}

For comparison, Mac Lane--Moerdijk explicitly separate the internal
Boolean condition, two-valuedness, and the splitting of all epimorphisms
\cite[pp.~270--276]{MacLaneMoerdijk1992}. Their later choice-independence
construction is not a substitute for the source's stronger Boolean
package \cite[pp.~291--295]{MacLaneMoerdijk1992}; the proof of that later
construction continues beyond this selected reading. The correspondence
above is relative to the given abstract-set base, not a classification
of all elementary topoi or of ordinary set-theoretic models.
The full-datum closed functor and nondegenerate two-valued collapse
follow in Theorems~\ref{thm:p10-lawvere-closed-base-partials}--
\ref{thm:p10-lawvere-closed-base-collapse}, after this BVM construction.
The countable-chain-condition/CH transfer, global-spectrum universal
property, and independently owed recursion/free-algebra dependencies
remain separate; neither choice nor the correspondence proves them.
